\section{A cubic obstruction on a projective threefold}\label{sec:compact-example}

Let $Q$ be the lattice hexagon with vertices
\[
 (0,-1),\ (1,0),\ (0,1),\ (-1,1),\ (-2,0),\ (-1,-1)
\]
in cyclic order. In $\Z^3$ take the pyramid
$P=\operatorname{conv}((Q,1),a)$ with apex $a=(0,0,-1)$.
Subdivide the face fan of $P$ at the ray through $m=(-1,0,0)$, and
let $Z$ be the resulting toric threefold.

\begin{proposition}\label{prop:compact-cubic}
A general anticanonical hypersurface $X\subset Z\times\mathbf P^1$
is a projective Calabi--Yau threefold with two singular points, both
of type $X_Q$. For $h=h^1(X,\Theta_X)$ its semiuniversal ring is
\[
 \OO_{\Def(X),0}\simeq
 \C\{w_1,\ldots,w_h,u,v,s\}/(us,vs,s^3).
\]
In particular there is a first-order deformation that extends over
$\C[\varepsilon]/(\varepsilon^3)$ and has no extension over
$\C[\varepsilon]/(\varepsilon^4)$.
\end{proposition}

\begin{proof}
The plane of the upper facet has equation $z=1$. In the order of
the edges of $Q$ the six side facets have equations
$\langle n_i,\,\cdot\,\rangle=1$ with outward normals
\[
 (2,-2,-1),\ (2,2,-1),\ (0,2,-1),\
 (-1,1,-1),\ (-1,-1,-1),\ (0,-2,-1).
\]
Thus $P$ is reflexive. The absolute determinants of the three rays
of its side cones are $1,1,1,2,2,1$, respectively. The two cones of
determinant two share the face spanned by $a$ and $(-2,0,1)$.
Its midpoint $m$ divides each into two unimodular cones. This star
subdivision is projective and crepant. The remaining maximal cone
$\sigma$ over $(Q,1)$ is smooth off its closed orbit, since the
edges of $Q$ are primitive. Write $p$ for its fixed point in $Z$.

The anticanonical bundle on $Z$ is the pullback of the globally
generated anticanonical bundle for the face fan of $P$. A general
section on $Z\times\mathbf P^1$ cuts out a smooth hypersurface away
from $\{p\}\times\mathbf P^1$ by Bertini. Its restriction there is
a general binary quadratic, with two simple roots in $\C^*$.
At each root projection to $Z$ identifies the hypersurface germ
with $(X_Q,0)$ by the implicit function theorem. In particular $X$
is normal. Adjunction gives $\omega_X\simeq\OO_X$. On the projective
Gorenstein toric fourfold $A=Z\times\mathbf P^1$ one has
$H^i(A,\OO_A)=0$ for $i>0$. Serre duality and the sequence
$0\to\omega_A\to\OO_A\to\OO_X\to0$ give
$H^0(X,\OO_X)=\C$ and $H^1(X,\OO_X)=0$.

To resolve the two points, subdivide $\sigma$ successively at
$(0,0,1)$ and $(-1,0,1)$. Their projections are the two interior
lattice points of $Q$. The first star joins the origin of the
polygon to its six edges. The second bisects the common edge of the
two triangles of normalized area two. All resulting triangles are
unimodular. Let $\widehat Z$ be the smooth projective toric
threefold so obtained and $Y\subset\widehat Z\times\mathbf P^1$
the proper transform of $X$. Near the two roots the resolution
is identified by projection with the toric resolution of $X_Q$.
Every interior edge in the triangulation contains an interior vertex
as an endpoint. Thus each exceptional fibre is the union of the two
toric divisors of the interior rays, meeting transversely along their
common invariant curve. In particular $Y\to X$ is a projective log
resolution and is an isomorphism on the regular locus.

Use the two projection charts to identify both oriented links with
the link $L$ of $X_Q$. We claim that the image of
\[
 H^2(Y,\C)\longrightarrow H^2(L,\C)\oplus H^2(L,\C)
\]
is the diagonal. Indeed the fan polytope of the original fourfold is
the convex hull of $(P,0)$ and $(0,0,0,\pm1)$. Its boundary lattice
points are the two latter vertices and the points of $(\partial P,0)$
in the lattice. The fan of $\widehat Z\times\mathbf P^1$ gives a
projective unimodular triangulation of this boundary using all its
lattice points. We may choose our general section nondegenerate with
respect to the polar Newton polytope. By the proof of
\cite[Theorem~4.4.2]{Bat}, the components of the intersections of $Y$
with the toric boundary divisors span $H^2(Y,\C)$. The six boundary
divisors corresponding to vertices of $Q$ have irreducible
intersections with $Y$, since their rays pass through vertices of the
fan polytope. They restrict in both projection charts to the same
six toric divisor classes on the punctured cone. Divisors associated
to the two interior points of $Q$ split over the two roots, and
each component is exceptional, hence disjoint from both links. All
the other boundary divisors miss the two germs. This proves
containment in the diagonal.

For the reverse inclusion, the fan of the punctured cone consists
of the proper faces of $\sigma$. Its homogeneous coordinate presentation
\cite[Theorem~2.1]{Cox} is a principal torus bundle over the punctured
cone. Indeed the fan is smooth and
its six rays span $\Z^3$, as the rays over $(1,0),(0,1),(-1,1)$
already form a lattice basis. The total space is the complement in
$\C^6$ of coordinate subspaces of complex codimension at least two
and has vanishing first and second cohomology by general position.
The low-degree sequence of the bundle then gives a surjection from
the complexified character group of the acting torus to $H^2(L,\C)$
by first Chern classes, using the retraction of the punctured cone
onto its link. The six coordinate characters generate this group
and give precisely the six ray divisor classes. This proves the claim.

By duality the relation space $K$ in \eqref{eq:K} is the
anti-diagonal in $H_2(L,\C)\oplus H_2(L,\C)$. Taking the same
marking and homogeneous toric form in the two projection charts
identifies their period images with the same germ
$G_Q=(\PD\circ\ell_Q)(B_Q)$. The restriction of a global volume form
differs from the toric generator by a unit in each chart and hence
has the same image by Theorem~\ref{thm:local}. Since the equations of $G_Q$ are
homogeneous, the pullbacks of its two ideals to $K$ coincide.
Theorem~\ref{thm:global} therefore gives $B_Q$ as the factor
transverse to the smooth germ of dimension $h$.

It remains to compute $B_Q$. The ordered edge vectors of $Q$ are
\[
 (1,1),\ (-1,1),\ (-1,0),\ (-1,-1),\ (1,-1),\ (1,0).
\]
Take the translation slice $t_6=0$. The general solution of the
first moment equation is
\[
 (t_1,\ldots,t_6)=(u+3s,v+s,4s,u+s,v+3s,0).
\]
The second moment is $(4s(u+v),4s(u-v))$. Modulo $(us,vs)$
the third is $(-12s^3,0)$. For every $k\geq 4$, setting $s=0$
cancels the two pairs with parameters $u,v$. Every other term of
the $k$th moment is divisible by $us$, $vs$ or $s^3$. This gives
the stated ring. Its quotient by $(w_1,\ldots,w_h,u,v)$ is
$\C\{s\}/(s^3)$, so $s^2\ne0$. Setting $s=\varepsilon$ and the
other parameters to zero gives a deformation modulo $\varepsilon^3$.
Every lift of its first-order class would send $s$ to
$\varepsilon+O(\varepsilon^2)$, whose cube is nonzero modulo
$\varepsilon^4$. This proves the assertion.
\end{proof}
